% F-02 四问递进关系图
\documentclass[border=5pt]{article}
\usepackage[margin=0pt,paperwidth=19cm,paperheight=11cm]{geometry}
\pagestyle{empty}
\usepackage{fontspec}
\usepackage{ctex}
\newcommand{\fontdir}{../../../00_知识库/论文写作/LaTeX模板_2026国赛版/fonts/}
\setCJKfamilyfont{song}[
  Path = \fontdir,
  UprightFont = SourceHanSerifCN-Regular.otf,
  BoldFont   = SourceHanSerifCN-Bold.otf,
  AutoFakeBold = true,
]{SourceHanSerifCN}
\newcommand*{\song}{\CJKfamily{song}}

\usepackage{tikz}
\usetikzlibrary{arrows.meta,positioning,calc}

\definecolor{primary}{HTML}{1F3A93}
\definecolor{primaryDark}{HTML}{1A3A6E}
\definecolor{secondary}{HTML}{D9E4F2}
\definecolor{accent}{HTML}{D35400}
\definecolor{accentLight}{HTML}{FDEBD0}
\definecolor{gray}{HTML}{95A5A6}
\definecolor{textDark}{HTML}{333333}
\definecolor{arrowColor}{HTML}{5D6D7E}
\definecolor{groupBg}{HTML}{F4F6F7}

\tikzset{
  process/.style = {rectangle,draw=primaryDark,line width=0.8pt,fill=primary,
    text=white,minimum width=3.0cm,minimum height=0.9cm,
    align=center,font=\song\footnotesize,inner sep=4pt,rounded corners=2pt},
  subprocess/.style = {rectangle,draw=primaryDark,line width=0.7pt,fill=secondary,
    text=textDark,minimum width=3.0cm,minimum height=0.85cm,
    align=center,font=\song\footnotesize,inner sep=4pt,rounded corners=2pt},
  startstop/.style = {rectangle,draw=accent,line width=1.0pt,fill=accentLight,
    text=accent!85!black,minimum width=3.0cm,minimum height=0.9cm,
    align=center,font=\song\footnotesize,inner sep=4pt,rounded corners=6pt},
  arrow/.style = {-{Stealth[length=2.2mm,width=2.0mm]},line width=0.9pt,draw=arrowColor},
  arrowBold/.style = {-{Stealth[length=2.5mm,width=2.2mm]},line width=1.2pt,draw=accent},
}

\begin{document}
\begin{tikzpicture}

  % === 第一行：四个问题盒（process风格） ===
  \node[process,minimum width=3.2cm,minimum height=2.0cm] (Q1) at (1.6, 0) {
    \textbf{Q1 预热阶段}\\[2pt]
    0--1800 s\\[2pt]
    附录2 定物性\\固定半径 $R_0{=}2$ cm\\
    $\rho,c_p,k$ 常数，$D{=}D(C)$
  };
  \node[process,minimum width=3.2cm,minimum height=2.0cm] (Q2) at (6.0, 0) {
    \textbf{Q2 变物性}\\[2pt]
    0--3 h\\[2pt]
    附录3 全变物性\\
    $\rho(C),c_p(C),$\\
    $k(C),D(C,T)$
  };
  \node[process,minimum width=3.2cm,minimum height=2.0cm] (Q3) at (10.4,0) {
    \textbf{Q3 长时程}\\[2pt]
    至达标\\[2pt]
    $\max C(r,t){<}0.15$\\
    事件驱动终止\\
    确定烘干时间
  };
  \node[process,minimum width=3.2cm,minimum height=2.0cm] (Q4) at (14.8,0) {
    \textbf{Q4 移动边界}\\[2pt]
    考虑收缩\\[2pt]
    附录4 + 附件2\\
    $R{=}R(t)$, 坐标变换
  };

  \draw[arrowBold] (Q1.east) -- (Q2.west);
  \draw[arrowBold] (Q2.east) -- (Q3.west);
  \draw[arrowBold] (Q3.east) -- (Q4.west);

  % === 第二行：灰色升级盒（subprocess风格，充足间距） ===
  \node[subprocess,minimum width=3.2cm] (G1) at (1.6, -3.0)
    {基础：建立 PDE 框架\\验证数值方法};
  \node[subprocess,minimum width=3.2cm] (G2) at (6.0, -3.0)
    {扩展：引入变物性\\IMEX 交替推进};
  \node[subprocess,minimum width=3.2cm] (G3) at (10.4,-3.0)
    {应用：全过程模拟\\灵敏度与不确定度};
  \node[subprocess,minimum width=3.2cm] (G4) at (14.8,-3.0)
    {修正：移动边界\\收缩效应量化};

  \draw[arrow] (Q1.south) -- (G1.north);
  \draw[arrow] (Q2.south) -- (G2.north);
  \draw[arrow] (Q3.south) -- (G3.north);
  \draw[arrow] (Q4.south) -- (G4.north);

  % === 第三行：统一框架（startstop风格，居中收窄） ===
  \node[startstop,minimum width=15cm,minimum height=1.0cm,
        fill=accentLight!60,draw=accent!50,text=textDark]
    (Uni) at (8.2, -5.5) {
    \textbf{统一框架：}~~~轴对称圆柱 1D 径向 ($L/R{=}12.5$)~~~
    热传导 + Fick 扩散 PDE 耦合\\[2pt]
    全隐式 FVM + Picard 迭代~~~同一模型、逐问递进、逐级加码
  };

  \draw[arrow] (G1.south) -- ++(0,-1.1) -| (Uni.north);
  \draw[arrow] (G2.south) -- ++(0,-1.1) -| (Uni.north);
  \draw[arrow] (G3.south) -- ++(0,-1.1) -| (Uni.north);
  \draw[arrow] (G4.south) -- ++(0,-1.1) -| (Uni.north);

\end{tikzpicture}
\end{document}